\documentclass[a4paper,12pt]{article}

\usepackage[slovene]{babel}
\usepackage{amsfonts, amssymb, amsmath, amsthm}
\usepackage[utf8]{inputenc}
\usepackage[T1]{fontenc}
\usepackage{lmodern}
\usepackage{graphicx}
\usepackage{xcolor}

% bibliography
\usepackage[
    backend=biber,
    sorting=nty,
]{biblatex}
\addbibresource{references.bib}


\usepackage{bbold}
% default sets
\newcommand{\N}{\mathbb{N}}
\newcommand{\Z}{\mathbb{Z}}
\newcommand{\Q}{\mathbb{Q}}
\newcommand{\R}{\mathbb{R}}
\newcommand{\C}{\mathbb{C}}
\newcommand{\F}{\mathbb{F}}
\newcommand{\HH}{\mathbb{H}}
\newcommand{\OO}{\mathcal{O}}

%%% Množice
\newcommand{\set}[1]{\left\{#1\right\}}
\newcommand{\setb}[2]{\set{#1 \ | \  #2}}

\newcommand{\eps}{\varepsilon}
% Implikacija
\newcommand{\lthen}{\implies}
\newcommand{\liff}{\iff}

%%% Quantifiers
\newcommand{\all}[1]{\forall #1 \,.\,}
\newcommand{\some}[1]{\exists #1 \,.\,}
\newcommand{\exactlyone}[1]{\exists! #1 \,.\,}
\DeclareMathOperator{\Fix}{Fix}
\newcommand{\one}{\mathbb{1}}

\def\qed{$\hfill\Box$}   % konec dokaza
\def\qedm{\qquad\Box}   % konec dokaza v matematičnem načinu
\newtheorem{izrek}{Izrek}[section]
\newtheorem{trditev}[izrek]{Trditev}
\newtheorem{posledica}[izrek]{Posledica}
\newtheorem{lema}[izrek]{Lema}
\newtheorem{opomba}[izrek]{Opomba}
\newtheorem{definicija}[izrek]{Definicija}
\newtheorem{zgled}[izrek]{Zgled}

\title{Izračun negibnih točk vektorja odsekoma linearnih funkcij \\ 
\Large Osnutek dela diplomskega seminarja}
\author{Ruslan Urazbakhtin \\
Mentor: prof.\ dr.\ Alexander Keith Simpson \\
Fakulteta za matematiko in fiziko \\
Oddelek za matematiko}
\date{}

\begin{document}


%%%%%%%%%%%%%%%%%%%%%%%%%%%%%%%%%%%%%%%%%%%%%%%%%%%%%%%%%%%%%%%%%%%%%


\maketitle


%%%%%%%%%%%%%%%%%%%%%%%%%%%%%%%%%%%%%%%%%%%%%%%%%%%%%%%%%%%%%%%%%%%%%
\section{Odsekoma linearne funkcije}
Osnovni pojem, s katerim se bomo ukvarjali v diplomski nalogi, so \emph{odsekoma linearne funkcije}. Preden jih formalno definiramo, si oglejmo preprosta zgleda.

\begin{zgled}
  \label{zgl:olf-simple}
  Naj bo funkcija \(f \colon \R \to \R\) podana po kosih:
  \[
    f(x) = \begin{cases}
      x; &x<0;\\
      x+1; &x \geq 0.
    \end{cases}
  \]
  Vidimo, da je funkcija \(f\) nezvezna v točki \(0\), medtem ko je na posamez\-nih intervalih podana z \emph{linearnim izrazom} oblike \(q_1 x + q_0\).
\end{zgled}

Bolj splošno lahko obravnavamo funkcije \(n\) spremenljivk \(f \colon \R^n \to \R\).

\begin{zgled}
  \label{zgl:olf-ndim}
  Naj bo funkcija \(f \colon \R^n \to \R\) podana s predpisom
  \[
    f(x_1, \ldots, x_n) = \left|\setb{i \in [n]}{x_i \geq 0}\right|.
  \]
  Domeno \(\R^n\) razbijemo glede na predznake koordinat.
  Za vsak izbor predznakov \((\varepsilon_1, \ldots, \varepsilon_n)\), kjer je \(\varepsilon_i \in \{-1,1\}\), dobimo množico oblike
    \[
      \{(x_1,\ldots,x_n)\in\R^n \mid (\eps_i \geq 0 \lthen x_i \geq 0) \land (\eps_i < 0 \lthen x_i < 0),\ i \in [n]\}.
    \]
  Te množice so ortanti v \(\R^n\) in skupaj tvorijo razbitje domene na \(2^n\) delov. Na posameznem ortantu funkcija \(f \) je podana s linearnim izrazom oblike
  \[
    q_n x_n + \cdots + q_1 x_1 + q_0,
  \]
  kjer velja \(q_n = \cdots = q_1 = 0\), medtem ko je konstanta \(q_0\) enaka številu indeksov \(i\), za katere je \(x_i \geq 0\). Očitno je, da je funkcija \(f\) nezvezna. 
\end{zgled}



Sedaj podajmo še formalno definicijo, ki jo povzemamo po \cite{zbMATH06805750}.
\begin{definicija}
  \emph{Linearni izraz} v spremenljivkah \(x_1, \ldots, x_n\) je izraz
  \[
    q_nx_n + \ldots + q_1 x_1 + q_0,
  \]
kjer so \(q_0, q_1 \ldots, q_n \in \R\). \emph{Pogojni linearni izraz} je par \(C \vdash e\), kjer je \(e\) linearni izraz in \(C\) končna množica neenačb med linearni izrazi, torej vsak element množice \(C\) je ene izmed oblik
    \[
        e_1 \leq e_2, \quad e_1 < e_2.
    \]
\end{definicija}

Za \(\vec{r} = (r_1, \ldots, r_n) \in \R^n\) z \(C(\vec{r})\) označimo konjunkcijo neenačb, ki nastane iz \(C\) z nadomestitvijo spremenljivk z vrednostmi \(r_1, \ldots, r_n\).

\begin{opomba}
  Z pomočjo pogojnega linearnega izraza \(C \vdash e\) lahko opišemo en kos funkcije:
  \begin{itemize}
      \item Domena kosa je \(\setb{\vec{r} \in \R^n}{C(\vec{r}) = \top}\).
      \item Linearni izraz \(e\) podaja linearno funkcijo nad domeno.
  \end{itemize}
\end{opomba}

\begin{definicija}
  \label{def:odsekoma-linearna-funkcija}
    Pravimo, da je funkcija \(f \colon D \subseteq \R^n \to \R\) \emph{odsekoma linearna}, če obstaja končna množica (oz.\ sistem) \(\mathcal{F}\) pogojnih linearnih izrazov v spremenljivkah \(x_1, \ldots, x_n\), da velja:
    \begin{enumerate}
        \item \(\all{\vec{r} \in D} \some{(C \vdash e) \in \mathcal F} C(\vec{r}) = \top\) in
        \item \(\all{\vec{r} \in D} \all{(C \vdash e) \in \mathcal F} C(\vec{r}) = \top \lthen f(\vec{r}) = e(\vec{r})\).
    \end{enumerate}
    Pravimo tudi, da \(\mathcal{F}\) \emph{predstavlja} funkcijo \(f\).
\end{definicija}

\begin{zgled}
  V smislu definicije \ref{def:odsekoma-linearna-funkcija} sta funkciji, podani v zgledih \ref{zgl:olf-simple} in \ref{zgl:olf-ndim}, očitno odsekoma linearni. Izraz "`odsekoma"' se nanaša na dejstvo, da je funkcija v splošnem nezvezna.
\end{zgled}

V nadaljevanju se bomo ukvarjali z algoritmi, ki računajo negibne točke odseko\-ma linearnih funkcij. Ti algoritmi morajo poznati ne le vrednost funkcije v dani točki, temveč tudi pogojni linearni izraz, s katerim lahko opišemo okolico te točke. Zato je smiselno razmisliti, kako lahko odsekoma linearne funkcije učinkovito predstavimo v kodi.

\subsection{Lokalni algoritem}
Ponovno se vrnimo k funkciji iz zgleda \ref{zgl:olf-ndim}. Izkazuje se, da je domena, podana s pogojnim linearnim izrazom \(C \vdash e\), konveksna množica. Zato za opis funkcije po definiciji \ref{def:odsekoma-linearna-funkcija} potrebujemo vsaj \(2^n\) ortantov. To pomeni, da bo pri predstavitvi odsekoma linearne funkcije s sistemom \(\mathcal{F}\) operacija iskanja okolice dane točke zahtevala \(\mathcal{O}(2^n)\) časa. 

Zato je smiselno vpeljati pojem \emph{lokalnega algoritma}, ki je bil implicitno podan v \cite{zbMATH06805750} in eksplicitno predstavljen v \cite{Petkovic2017}.

\begin{definicija}
    \emph{Lokalni algoritem} za odsekoma linearno funkcijo \(f \colon D \subseteq \R^{n} \to \R\), je algoritem, ki za podano točko \(\vec{r} \in D\) vrača pogojni linearni izraz \(C_{\vec{r}} \vdash e_{\vec{r}}\), da velja
    \begin{enumerate}
        \item \(C_{\vec{r}}(\vec{r}) = \top\);
        \item \(\all{\vec{s} \in D} C_{\vec{r}}(\vec{s}) = \top \lthen f(\vec{s}) = e_{\vec{r}}(\vec{s})\).
    \end{enumerate}
    Pri tem je množica \(\setb{C_{\vec{r}} \vdash e_{\vec{r}}}{\vec{r} \in D}\) končna.
\end{definicija}

\begin{zgled}
  Naj bo funkcija \(f\) ista kot v zgledu \ref{zgl:olf-ndim}. Lokalni algoritem za funkcijo \(f\) definiramo na naslednji način. Naj bo \(\vec{r} \in \R^n\). Definiramo pogojni linearni izraz \(C_{\vec{r}} \vdash e_{\vec{r}}\) kot
  \[
    C_{\vec{r}} := \set{x_i \geq 0}_{r_i \geq 0} \cup \set{x_i < 0}_{r_i < 0}.
  \]
  in
  \[
    e_{\vec{r}} := \left|\setb{i \in [n]}{r_i \geq 0}\right|.
  \]
  %
  Jasno je, da je število različnih pogojnih linearnih izrazov oblike \(C_{\vec{r}} \vdash e_{\vec{r}}\), ko \(\vec{r}\) preteče \(\R^n\), enako \(2^n\). S tem smo dejansko podali lokalni algoritem.

  Prednost takšne predstavitve funkcije \(f\) je v tem, da operacija iskanja okolice dane točke sedaj zahteva le \(\mathcal{O}(n)\) časa.
\end{zgled}

Omenimo še, da lahko vsako eksplicitno predstavitev funkcije \(f\) s sistemom \(\mathcal{F}\) enostavno pretvorimo v lokalni algoritem. Obratna pretvorba pa je v splošnem zahtevna.


\section{Polne mreže, negibne točke in izrek Knaster-Tarskega}

Naslednje vprašanje je, kakšni so zadostni pogoji za odsekoma linearno funkcijo, da ima negibno točko. Nedvoumno je, da to ne drži za vse odsekoma linearne preslikave. Preprost protiprimer je funkcija 
\[
f \colon \R \to \R,\quad f(x) = x + 1.
\]
%
Obenem znana izreka o negibnih točkah, kot sta Banachovo skrčitveno načelo in Brouwerjev izrek o negibni točki, tu ne pomagata, saj so odsekoma linearne funkcije v splošnem nezvezne. Odgovor na naše vprašanje bomo našli z uporabo teorije urejenosti.

Preden formuliramo osnovni izrek, navedemo osnovni definiciji. Vsebino tega razdelka povzemamo po \cite[1.\ poglavje]{zbMATH01576861}.

\begin{definicija}
    Naj bo \((E, \leq)\) delno urejena množica:
    \begin{itemize}
        \item \((E, \leq)\) je \emph{mreža}, če za poljubna elementa \(x, y \in E\), ima množica \(\set{x, y}\) največjo spodnjo mejo \(x \land y\) in najmanjšo zgornjo mejo \(x \lor y\). 
        \item \((E, \leq)\) je \emph{polna mreža}, če vsaka podmnožica \(X \subseteq E\) ima  največjo spodnjo mejo \(\bigwedge X\) in najmanjšo zgornjo mejo \(\bigvee X\). 
    \end{itemize}
\end{definicija}

\begin{definicija}
    Naj bosta \((E, \leq_E)\) in \((F, \leq_F)\) delno urejeni množici. Preslikava \(f \colon E \to F\) je \emph{monotona}, če
    \[
        \all{x, y \in E} x \leq_E y \lthen f(x) \leq_F f(y).
    \]
\end{definicija}

\begin{opomba}
  Beseda "`monotonost"' v tem delu ima enak pomen kot beseda "`naraš\-čajočost"'.
\end{opomba}

Sedaj lahko podamo osnovni izrek, ki bo dal odgovor na zastavljeno vprašanje.
\begin{izrek}[Knaster-Tarski]
  \label{izr:knaster-tarski}
  Naj bo \((E, \leq)\) polna mreža in \(f \colon E \to E\) monotona preslikava. Tedaj \(\bigwedge \Fix(f)\) in  \(\bigvee \Fix(f)\) sta elementi \(\Fix(f)\), kjer je 
  \[
    \Fix(f) = \setb{x \in E}{f(x) = x}.
  \]
\end{izrek}

Vidimo, da bo za monotono odsekoma linearno funkcijo 
\[
  f \colon E \subseteq \R \to E,
\]
kjer je \(E\) polna mreža, množica \(\Fix(f)\) neprazna in bo vsebovala minimum ter maksimum (torej najmanjšo in največjo negibno točko).  

Ustrezna izbira za množico \(E\) je na primer zaprt interval \([0,1]\), ki je očitno polna mreža, saj \(\R\) izpolnjuje Dedekindov aksiom.

Ko vemo, za kateri odsekoma linearni funkciji lahko izračunamo njene negibne točke, lahko problem formuliramo tudi v algoritmični obliki.
\section{Algoritmi}
V tem razdelku bomo vzpostavili algoritmični problem izračuna negibnih točk odse\-koma linearnih funkcij in ga rešili z ustreznimi algoritmi, katerih koda bo tudi javno dostopna. Zanimali nas bosta le najmanjša in največja negibna točka, saj sta pomembni za relevantne probleme, v katerih želimo učinkovito izračunati ti vrednosti. Pojma najmanjše in največje negibne točke sta hkrati dualna: če imamo algoritem za izračun najmanjše negibne točke, lahko s pomočjo njega in ustreznih manipulacij s preslikavo izračunamo tudi največjo negibno točko. Obenem bomo zastavili vprašanja, na katera odgovor še ni znan, in poskušali poiskati ustrezne rešitve.  

\subsection{Formulacija problema}
Naj bo \(E\) polna mreža. Na množici monotonih preslikav v drugi spremenljivki \(E^{X \times E}\) definiramo operator 
\[
  \dagger_c: E^{X \times E} \to E^X, \quad f \mapsto f^{\dagger_c},
\] 
kjer je 
\[
  \all{x \in X} f^{\dagger_c}(x) = c(\Fix(y \mapsto f(x, y) \colon E \to E))
\]
in je \(c\) neka funkcija izbire elementa iz množice.

\begin{zgled}
  Naj bo \(c = \bigwedge\). Tedaj operator \(\dagger_c\) za dano monotono preslikavo v drugi spremenljivki \[f \colon X \times E \to E\] 
  vrne preslikavo 
  \[
    f^{\dagger_c} \colon X \to E,
  \] 
  ki za vsak \(x \in X\) izračuna najmanjšo negibno točko preslikave 
  \[
    f_x \colon E \to E, \quad f_x(y) = f(x, y). \qedhere
  \]
\end{zgled}

Zanimali nas bodo algoritmi, ki predstavljajo operator \(\dagger_c\) na množici odsekoma linearnih funkcij, monotonih v drugi spremenljivki (oziroma vektorjev takih funkcij), \(E^{X \times E}\), za različne izbire množic \(X\) in \(E\) ter funkcije \(c\). Oglejmo si nekaj možnih izbir.

Najenostavnejša izbira je \(X = \one\) in \(E = [0,1]\). Tako dobimo množico funkcij
\[
  [0,1]^{\one \times [0,1]}.
\]
Operator \(\dagger_c\) tedaj vrne funkcijo na množici \(\one\), pri čemer lahko vrednost \(f^{\dagger_c}\) razumemo kot negibno točko (izbrano s pomočjo funkcije \(c\)) funkcije
\[
  f \colon [0,1] \to [0,1].
\]

Prejšnji primer lahko z izbiro \(X = [0,1]^n\) posplošimo na funkcije več spremenljivk
\[
  f \colon [0,1]^n \times [0,1] \to [0,1].
\]
Operator \(\dagger_c\) tedaj vrne funkcijo na množici \([0,1]^n\):
\[
  f^{\dagger_c} \colon [0,1]^n \to [0,1].
\]
Izkaže se, da je tudi funkcija \(f^{\dagger_c}\) odsekoma linearna. Ob dodatni predpostavki, da je monotona v zadnji spremenljivki, lahko ponovno uporabimo operator \(\dagger_c\) (ki zdaj deluje na množici \([0,1]^{[0,1]^n}\)) in tako dobimo funkcijo
\[
  f^{\dagger_c \dagger_c} \colon [0,1]^{n-1} \to [0,1].
\]

Ob ustreznih predpostavkah o monotonosti lahko ta postopek nadaljujemo. Izka\-že se, da je to smiselno, vendar algoritem, predstavljen v \cite{zbMATH06805750} in posodobljen v \cite{Petkovic2017}, povzroča eksponentno časovno zahtevnost ter pri večjem številu spremenljivk hitro postane neuporaben.

Zato bomo problem poskusili rešiti z izbiro množic \(X = [0,1]^m\) in \(E = [0,1]^n\), oziroma ga bomo posplošili na vektor odsekoma linearnih funkcij
\[
  F = (f_1, \ldots, f_n) \colon [0,1]^m \times [0,1]^n \to [0,1]^n,
  \quad
  f_i \colon [0,1]^m \times [0,1]^n \to [0,1].
\]
Operator \(\dagger_c\) v tem primeru vrne preslikavo na množici \([0,1]^m\):
\[
  F^{\dagger_c} \colon [0,1]^m \to [0,1]^n.
\]

\subsection{Preizkusi časovne zahtevnosti}
Vse različice algoritma bomo testirali na skrbno izbranih primerih, zasnovanih tako, da poudarijo njihove časovne lastnosti. Pričakujemo, da bo posplošitev na vektorski primer prinesla občuten izboljšek zmogljivosti v primerjavi z iterativno uporabo skalarne različice.


\printbibliography

\end{document}

\section*{Angleško-slovenski slovar strokovnih izrazov}